% =====================================================================
%  Annexe D - Interfaces complementaires
% =====================================================================
%  Ecrans qui n'ont pas trouve place dans les chapitres de sprint, soit
%  parce qu'ils prolongent un ecran deja presente par defilement, soit
%  parce qu'ils relevent du confort d'usage plus que de la demonstration
%  d'une regle metier. Chaque capture occupe ici sa propre figure, pour
%  rester lisible a l'impression.
% =====================================================================
\cleardoublepage
\section*{Annexe D -- Interfaces complementaires}
\addcontentsline{toc}{section}{Annexe D -- Interfaces complementaires}
\serieAnnexe{D}

Les ecrans reunis dans cette annexe completent ceux presentes dans les
chapitres de sprint. Ils ont ete ecartes du corps du rapport non parce
qu'ils seraient secondaires dans l'usage, mais parce que la regle metier
qu'ils servent y etait deja demontree par un autre ecran.

\subsection*{Portail client}

La page d'accueil du portail, figure~\ref{fig:annexe-accueil}, est le
premier contact d'un client avec l'entreprise. Elle presente les
prestations, les zones desservies, et conduit au formulaire de
reservation. Sa refonte faisait partie du perimetre du stage : le site
anterieur ne comportait aucun point d'entree vers une demande de
transport.

\capture{s0-page-accueil.jpg}
  {Page d'accueil du portail client}
  {fig:annexe-accueil}

Le formulaire de reservation se poursuit au-dela de ce qui a ete montre
au chapitre~\ref{chap:sprint3}. La figure~\ref{fig:annexe-reservation2}
en presente la suite : caracteristiques de la marchandise, contraintes de
manutention et coordonnees du destinataire.

\capture{s3-portail-reservation-2.jpg}
  {Suite du formulaire de reservation : marchandise et destinataire}
  {fig:annexe-reservation2}

Avant l'envoi, le client relit un recapitulatif complet,
figure~\ref{fig:annexe-recap}. Cette etape a ete ajoutee apres les
premiers essais d'utilisation : plusieurs demandes de test comportaient
une ville d'arrivee erronee, choisie par inadvertance dans la liste
deroulante.

\capture{s3-portail-reservation-3.jpg}
  {Recapitulatif de la demande avant validation par le client}
  {fig:annexe-recap}

\subsection*{Supervision du suivi en direct}

Le chapitre~\ref{chap:sprint3} presente le suivi en direct tel que le
\essentiel{client} le voit. La figure~\ref{fig:annexe-chauffeur} en montre
le pendant cote \essentiel{exploitation} : le meme trajet, mais depuis le
back-office, avec la liste de toutes les missions actives en bandeau
superieur. L'exploitant passe de l'une a l'autre sans quitter la carte,
et dispose en marge du nom du chauffeur, de sa vitesse instantanee, du
kilometrage parcouru et de l'heure d'arrivee estimee.

La difference entre les deux ecrans n'est pas cosmetique : le client ne
voit que sa mission, l'exploitant les voit toutes. C'est la meme donnee de
position, filtree par le role de celui qui la consulte.

\capture{s3-app-chauffeur.jpg}
  {Suivi en direct cote back-office : supervision de toutes les missions
   actives}
  {fig:annexe-chauffeur}

\subsection*{Exploitation de la flotte}

Le suivi du carburant, figure~\ref{fig:annexe-carburant1}, enregistre
chaque plein avec sa date, son volume, son montant et son kilometrage.
C'est de ce releve que decoule la consommation moyenne par vehicule.

\capture{s2-carburant-1.jpg}
  {Suivi du carburant : releve des pleins}
  {fig:annexe-carburant1}

Les statistiques de consommation, figure~\ref{fig:annexe-carburant2},
mettent en regard les vehicules du parc. Un ecart durable par rapport a la
moyenne d'une meme categorie constitue un indice de derive mecanique, et
alimente indirectement la maintenance preventive.

\capture{s2-carburant-2.jpg}
  {Statistiques de consommation comparees de la flotte}
  {fig:annexe-carburant2}

Le troisieme onglet de la fiche vehicule, figure~\ref{fig:annexe-fiche3},
regroupe les assurances : compagnie, numero de contrat, garanties et date
de fin de validite. Ces dates alimentent le module d'alertes au meme titre
que celles des pieces administratives.

\capture{s2-fiche-vehicule-3.jpg}
  {Fiche d'un vehicule : assurances et dates de fin de validite}
  {fig:annexe-fiche3}

\subsection*{Traitement des demandes et pilotage}

La figure~\ref{fig:annexe-devis} montre l'ecran d'etablissement d'un devis
a partir d'une demande recue : le responsable y retient ou corrige le prix
estime, fixe la duree de validite et envoie la proposition au client.

\capture{s3-demandes-transport-2.jpg}
  {Etablissement d'un devis a partir d'une demande de transport}
  {fig:annexe-devis}

Les deux volets suivants du tableau de bord completent celui presente au
chapitre~\ref{chap:sprint6}. Le premier, figure~\ref{fig:annexe-bord2},
detaille l'activite de transport par periode et par client.

\capture{s6-tableau-bord-2.jpg}
  {Tableau de bord : activite de transport par periode}
  {fig:annexe-bord2}

Le second, figure~\ref{fig:annexe-bord3}, rassemble les indicateurs de
maintenance et de facturation : cout des interventions, vehicules
immobilises, encours de reglement.

\capture{s6-tableau-bord-3.jpg}
  {Tableau de bord : indicateurs de maintenance et de facturation}
  {fig:annexe-bord3}
